\documentclass[preview]{standalone}

\usepackage{amsmath}
\usepackage{amssymb}
\usepackage{stellar}
\usepackage{definitions}

\begin{document}

\id{psicologia-personalita-tratti}
\genpage

\section{Teorie dei tratti}

\begin{snippetdefinition}{personalita-definition}{Personalità}
    La \emph{personalità} è l'insieme complessivo delle caratteristiche
    psicologiche che influenzano i pattern comportamentali tipici di un
    individuo nel corso del tempo e in situazioni differenti.
\end{snippetdefinition}

\begin{snippet}{teorie-personalita}
    Le teorie della personalità sono ipotesi sulla struttura e sul funzionamento
    delle personalità individuali. Cercano di comprendere l'unicità di ciascun
    individuo e di spiegare in che modo una personalità produca pattern
    comportamentali caratteristici.
\end{snippet}

\begin{snippetdefinition}{tratto-personalita-definition}{Tratto}
    Un \emph{tratto} è una caratteristica o attributo duraturo che influenza il
    comportamento in molte situazioni differenti.
\end{snippetdefinition}

\begin{snippetdefinition}{variabile-interveniente-definition}{Variabile interveniente}
    Una \emph{variabile interveniente} è una variabile non direttamente
    osservabile che permette di spiegare il legame tra stimoli e risposte
    osservabili.
\end{snippetdefinition}

\begin{snippetexample}{timidezza-tratto-example}{Timidezza}
    Il tratto della timidezza può spiegare perché una persona risponda in modo
    simile in situazioni sociali diverse, per esempio parlando poco, evitando il
    contatto visivo o mostrando disagio.
\end{snippetexample}

\begin{snippet}{allport-tratti}
    Gordon Allport considerava i tratti come i mattoni della personalità e come
    fonte dell'individualità. Distinse:
    \begin{itemize}
        \item \emph{tratti cardinali}, attorno ai quali una persona organizza
        la propria vita;
        \item \emph{tratti centrali}, cioè caratteristiche fondamentali come
        onestà o pessimismo;
        \item \emph{tratti secondari}, cioè caratteristiche più specifiche e
        meno decisive, come preferenze alimentari.
    \end{itemize}
\end{snippet}

\includesnpt[width=58\%,src=/snippet/static-psychology/allport-trait-types.jpg]{centered-img}

\begin{snippet}{cattell-eysenck}
    Raymond Cattell ridusse il vasto vocabolario dei tratti a \(16\) fattori
    fondamentali, chiamati \emph{tratti originari}. Hans Eysenck individuò
    invece tre macro-dimensioni: \emph{estroversione}, \emph{nevroticismo} e
    \emph{psicoticismo}.
\end{snippet}

\begin{snippetdefinition}{modello-cinque-fattori-definition}{Modello a cinque fattori}
    Il \emph{modello a cinque fattori}, o \emph{Big Five}, è un modello della
    personalità secondo cui le differenze individuali possono essere descritte
    attraverso cinque grandi dimensioni: apertura mentale, coscienziosità,
    estroversione, piacevolezza e nevroticismo.
\end{snippetdefinition}

\includesnpt[width=72\%,src=/snippet/static-psychology/big-five-personality-dimensions.jpg]{centered-img}

\begin{snippet}{big-five-significato}
    Il modello a cinque fattori tenta di organizzare il grande numero di tratti
    descrittivi in un sistema più semplice e generale. Ogni fattore comprende
    molti tratti specifici e presenta due poli opposti.
\end{snippet}

\begin{snippetexample}{estroversione-amigdala-example}{Estroversione e amigdala}
    In uno studio con risonanza magnetica funzionale, partecipanti con diversi
    livelli di estroversione osservavano volti impauriti, felici e neutri. I
    volti impauriti attivavano l'amigdala in modo simile per tutti i livelli di
    estroversione; i volti felici, invece, producevano maggiore attività
    dell'amigdala sinistra negli individui più estroversi. Il risultato suggerì
    un legame tra tratti di personalità e risposta neurale a stimoli sociali
    positivi.
\end{snippetexample}

\includesnpt[width=58\%,src=/snippet/static-psychology/extraversion-amygdala-activity.jpg]{centered-img}

\begin{snippet}{big-five-evoluzione}
    In prospettiva evoluzionistica, le dimensioni dei Big Five possono essere
    collegate a problemi sociali ricorrenti: essere di buona compagnia,
    collaborare, impegnarsi in sforzi prolungati, essere emotivamente affidabili
    e sviluppare idee. La variabilità individuale può riflettere l'adattamento a
    contesti ambientali diversi.
\end{snippet}

\begin{snippet}{valutazione-teorie-tratti}
    Le teorie dei tratti permettono descrizioni sintetiche delle differenze
    individuali, ma spiegano solo parzialmente come il comportamento venga
    generato o come la personalità si sviluppi. Descrivono caratteristiche
    associate ai comportamenti, ma non sempre chiariscono i processi che li
    producono.
\end{snippet}

\end{document}
